\section{Measured boot}

\subsection{Concept}

\begin{frame}{Measured boot}
  \begin{itemize}
    \item Another way of establishing \textbf{trust}
    \item Idea: measure caracteristics of the system and record the measurement
    \item One major target for measurement is the software being run
    \item Other targets include:
    \begin{itemize}
      \item configuration (e.g. U-Boot environment)
      \item boot parameters (e.g. kernel command line)
    \end{itemize}
    \item Also called \textbf{trusted} boot \code{!=} secure/verified boot
  \end{itemize}
\end{frame}

\begin{frame}{Differences from Secure Boot}
  \begin{itemize}
    \item The software is not \textbf{authenticated}, only \textbf{measured}
    \item To be useful, the measurement needs to be used to make some
          decision
    \item This is a different step called \textbf{attestation} or \textbf{appraisal}
    \item It can only consist of recording the measurement discrepancy, or more complex actions
  \end{itemize}
\end{frame}

\begin{frame}{Measurements}
  \begin{itemize}
    \item A measurement is a hash of the data we are trying to measure
    \item We need to use a cryptographically secure hash function
    \item In a later step, we'll want to compare this hash to a known-good value
    \item We need to store these hashes
  \end{itemize}
\end{frame}

\begin{frame}{Measured boot support}
  \begin{itemize}
    \item A lot of the software components we've already covered support Measured boot:
    \begin{itemize}
      \item TF-A
      \item U-Boot
    \end{itemize}
    \item systemd supports measured boot using
          \href{https://www.freedesktop.org/software/systemd/man/latest/systemd-measure.html}{systemd-measure}
    \begin{itemize}
      \item this feature is experimental
      \item currently only supports UKI, which are mostly used in UEFI boots
    \end{itemize}
  \end{itemize}
\end{frame}

\subsection{Trusted Platform Module (TPM)}

\begin{frame}{Trusted Platform Module}
  \begin{itemize}
    \item Dedicated and isolated crypto component
    \item Following the \href{https://trustedcomputinggroup.org/}{Trusted Computing Group} (TCG)'s specification:
    \begin{itemize}
        \item \href{https://trustedcomputinggroup.org/resource/tpm-main-specification/}{TPM 1.2}, originally published in 2003
        \item TPM 2.0 Library, also known as ISO/IEC 11889
    \end{itemize}
    \item In 1.2, the TPM was explicitly specified as a hardware component
          including a crypto co-processor among others.
    \item The current 2.0 spec explicitly drops this requirement\\
          % Here, ask "can you think of something we've seen that fits this description?"
          % This is from Part 0, Introduction, 4.3 Trusted Platform Module.
          % In version 184, this is on page 26.
          \code{Another reasonable implementation of a TPM is to have the code run on the host processor
                while the processor is in a special execution mode}
  \end{itemize}
\end{frame}

\begin{frame}{TPM 2.0}
  \begin{itemize}
    \item Current version of the specification
    \item This is e.g. the one that is a Windows 11 requirement
    \item Specifies the TPM as a \textit{library}
    \item The TCG's \href{https://trustedcomputinggroup.org/wp-content/uploads/2019_TCG_TPM2_BriefOverview_DR02web.pdf}{Brief overview}
          recognizes 5 types of TPM:
    \begin{itemize}
      \item Discrete TPMs, or dTPM: dedicated hardened chip\\
            example: the \href{<https://www.st.com/en/secure-mcus/st33ktpm2x.html}{ST33KTPM2X}
      \item Integrated TPMs: co-processor. Examples:
            % For some reason, it seems like both are being referenced to as an fTPM,
            % but it clearly seems to be silicon that is integrated in some Intel and Ryzen SoCs
      \begin{itemize}
        \item Microsoft's \href{https://learn.microsoft.com/en-us/windows/security/hardware-security/pluton/microsoft-pluton-security-processor}{Pluton}
        \item The AMD \href{https://dayzerosec.com/blog/2023/04/17/reversing-the-amd-secure-processor-psp.html}{Secure Processor} (ASP, or PSP)
      \end{itemize}
      \item Firmware TPMs, or fTPMs: hardware-isolated software. Examples:
      \begin{itemize}
        \item the \href{https://github.com/OP-TEE/optee_ftpm}{OP-TEE fTPM}
      \end{itemize}
      \item Software TPMs: only useful for testing\\
            example: this \href{https://github.com/microsoft/ms-tpm-20-ref}{reference implementation}
            from Microsoft.
      \item Virtual TPMs, or vTPMs: basically an fTPM running in a hypervisor\\
            Google \href{https://cloud.google.com/blog/products/identity-security/virtual-trusted-platform-module-for-shielded-vms-security-in-plaintext?hl=en}{implemented one}
            for GCP
    \end{itemize}
  \end{itemize}
\end{frame}

% Oops https://www.amd.com/en/resources/product-security/bulletin/amd-sb-4011.html
\begin{frame}{Using the TPM}
  \begin{itemize}
    \item Key storage and generation
    \item Platform Configuration Registers (PCRs)
    \begin{itemize}
      \item "Shielded Locations", meaning hardware protected by the TPM
            % This shows that this was clearly designed with measured boot in mind
      \item Meant to contain a log of measurements
      \item Some persist across reboot, most should be reset to initial value
      \item Initial value must be all 0s or all 1s, except for PCR[0], which
            can be used as a locality indicator
      \item Cannot be set, only \textbf{reset} or \textbf{extended}
        \begin{itemize}
          \item This means hashing the current value of the PCR appended with the measurement
          \item The order of measurement therefore impacts the final PCR value
        \end{itemize}
    \end{itemize}
  \end{itemize}
\end{frame}

\begin{frame}{Additional references}
  \begin{itemize}
    \item The Open Security Training 2 "Trusted Computing" track:
          \href{https://p.ost2.fyi/courses/course-v1:OpenSecurityTraining2+TC1101_IntroTPM+2024_v2/about}{TC1101},
          \href{https://p.ost2.fyi/courses/course-v1:OpenSecurityTraining2+TC1102_IntermediateTPM+2024_v1/about}{TC1102},
          \href{https://p.ost2.fyi/courses/course-v1:OpenSecurityTraining2+TC2202_tpm2-pytss+2025_v1/about}{TC2202}
    \item The Open Access \href{https://link.springer.com/book/10.1007/978-1-4302-6584-9}{Practical Guide to TPM 2.0}
  \end{itemize}
\end{frame}

\subsection{IMA/EVM}

\begin{frame}{Integrity Measurement Architecture (IMA)}
  \begin{itemize}
    \item \href{https://ima-doc.readthedocs.io/en/latest/index.html}{IMA} is a Linux
          kernel subsystem, found under \code{security/integrity/ima}
    \item Uses extended attributes (\kdochtml{filesystems/ext4/attributes.html}) to store
          measurement of files at runtime.
    \item Enabled via \kconfig{CONFIG_IMA}
    \item Continues the measurement process after the kernel takes over
    \item Focuses on the integrity of file \textbf{contents}
  \end{itemize}
\end{frame}

\begin{frame}{Integrity Measurement Architecture Appraisal}
  \begin{itemize}
    \item IMA appraisal validates the integrity of file content
    \item File content is verified either with:
      \begin{itemize}
        \item A hash: file content is protected against offline modifications
        \item A signature: file content is protected against both offline and
          online modifications
      \end{itemize}
    \item Only the file content is validated: nothing prevent from renaming
      a file or replacing it with another valid file
    \item IMA appraisal behaviour is controlled by the \texttt{ima\_appraise}
      kernel command line parameter:
      \begin{itemize}
        \item \texttt{enforce} appraisal is fully enabled: access to file with
          invalid or missing hashes of signatures is denied
        \item \texttt{log} access to file with invalid or missing hashes or
          signature is allowed, but logged
        \item \texttt{fix} hashes of files covered by the policy are updated,
          when a file is accessed
        \item \texttt{off}: appraisal is completely disabled: hashes or signature
          are neither generated nor validates
      \end{itemize}
  \end{itemize}
\end{frame}

\begin{frame}{Integrity Measurement Architecture Policies}
  \begin{itemize}
    \item Policies defines what is measured
    \item Controlled by the \texttt{ima\_policy} kernel command line parameter
    \item Several policies can be combined
    \item IMA comes with predefined policies:
      \begin{itemize}
        \item Controls what is measured
        \item \texttt{tcb}: measures all executed programs, files mmap'd for
          execution, and all files read with uid or effective uid set to 0
        \item \texttt{appraise\_tcb} appraises the integrity of all files owned by root
        \item \texttt{secure\_boot} appraises the integrity of files based on file signatures
        \item \texttt{fail\_securely} always force file signature verification
        \item \texttt{critical\_data} measures kernel integrity critical data
      \end{itemize}
    \item Custom policies can be defined
      \begin{itemize}
        \item Only well-known stable files should be measured
          \begin{itemize}
            \item Binaries, libraries, configuration\ldots
          \end{itemize}
        \item Files with frequent modifications should not be measured
          \begin{itemize}
            \item Logs, databases\ldots
          \end{itemize}
      \end{itemize}
  \end{itemize}
\end{frame}

\begin{frame}{Extended Verification Module (EVM)}
  \begin{itemize}
    \item \href{https://ima-doc.readthedocs.io/en/latest/ima-concepts.html\#extended-verification-module-evm}{EVM}
          is also a Linux kernel subsystem, found under \code{security/integrity/evm}
    \item Uses extended attributes (\kdochtml{filesystems/ext4/attributes.html}) to store
          measurement of files at runtime.
    \item Enabled via \kconfig{CONFIG_EVM}
    \item Continues the measurement process after the kernel takes over
    \item Focuses on the integrity of file \textbf{metadata}
  \end{itemize}
\end{frame}

\begin{frame}{Extended Verification Module (EVM)}
  \begin{itemize}
    \item Similarly to IMA, EVM is split into:
    \begin{itemize}
      \item EVM HMAC: allows to protect against offline modifications
      \item EVM Signature: also protects against runtime modifications
    \end{itemize}
  \item EVM can be enabled with the \texttt{securityfs} pseudo-filesystem:
    \texttt{/sys/kernel/security/evm}:
      \begin{itemize}
        \item Value is a bit mask of enabled configuration, bits can only
          transition from 0 to 1
        \item bit 0: Enable HMAC validation and creation
        \item bit 1: Enable digital signature validation
        \item bit 31: Disable further runtime modification of EVM policy
      \end{itemize}
    \item As for IMA, nothing prevent from renaming the file: filename is not
      part of the metadata
  \end{itemize}
\end{frame}

\begin{frame}{IMA/EVM vs dm-verity}
  \begin{itemize}
    \item Both aim to enforce userland integrity
    \item IMA/EVM enables:
    \begin{itemize}
      \item remote attestation
      \item auditing
    \end{itemize}
    \item This comes at the cost of a more complicated setup
    \item IMA policies must be crafted very carefully
    \item Neither is very useful without a full secure boot chain
  \end{itemize}
\end{frame}

